\QHeading{第15题　脱式计算。}
\par\noindent


\CalcBlock{1}{\begin{aligned}[t]&\phantom{={}}\color{black}36+40\times(97-57)\\&\color{answerred}=36+40\times40\\&\color{answerred}=36+1600\\&\color{answerred}=1636\end{aligned}}
\par

\CalcBlock[简便解法（补充）]{2}{\begin{aligned}[t]&\phantom{={}}\color{black}(864\times 4+168)\div 4\\&\color{answerred}=864\times 4\div 4+168\div 4\\&\color{answerred}=864+42\\&\color{answerred}=906\end{aligned}}
\MethodBlock{按运算顺序计算（基础练习）}{本单元还要练习运算顺序；与上面的简便解法对照。}{\begin{aligned}[t]&\phantom{={}}(864\times4+168)\div4\\&=(3456+168)\div4\\&=3624\div4\\&=906\end{aligned}}

\MethodBlock{方法二：先分别相除再合并}{}{\begin{aligned}[t]&\phantom{={}}(864\times 4+168)\div 4\\
&=864\times 4\div 4+168\div 4\\
&=3456\div 4+42\\
&=864+42\\
&=906\end{aligned}}

\CalcBlock{3}{\begin{aligned}[t]&\phantom{={}}\color{black}188\div(820-22\times33)\\&\color{answerred}=188\div(820-726)\\&\color{answerred}=188\div94\\&\color{answerred}=2\end{aligned}}
\CalcBlock[方法一]{4}{\begin{aligned}[t]&\phantom{={}}\color{black}525\div[(81-56)\times3]\\&\color{answerred}=525\div[25\times3]\\&\color{answerred}=525\div75\\&\color{answerred}=7\end{aligned}}
\MethodBlock{方法二：除以积，改成连续除}{连续除以两个数，等于除以这两个数的积。}{\begin{aligned}[t]&\phantom{={}}525\div [(81-56)\times 3]\\
&=525\div (81-56)\div 3\\
&=525\div 25\div 3\\
&=21\div 3\\
&=7\end{aligned}}

\CalcBlock[方法一]{5}{\begin{aligned}[t]&\phantom{={}}\color{black}[(472+78)\div50]\div11\\&\color{answerred}=[550\div50]\div11\\&\color{answerred}=11\div11\\&\color{answerred}=1\end{aligned}}
\MethodBlock{方法二：把连续除合成一次除法}{}{\begin{aligned}[t]&\phantom{={}}[(472+78)\div 50]\div 11\\
&=(472+78)\div (50\times 11)\\
&=(472+78)\div 550\\
&=550\div 550\\
&=1\end{aligned}}

\CalcBlock{6}{\begin{aligned}[t]&\phantom{={}}\color{black}12\times[164\div(12-8)]\\&\color{answerred}=12\times[164\div4]\\&\color{answerred}=12\times41\\&\color{answerred}=492\end{aligned}}
\par\noindent

\QHeading{第16题　在\(\bigcirc\)里填上“\(>\)”“\(<\)”或“\(=\)”。}

\Equation{36\times9+36\textcolor{answerred}{=}36\times(9+1)}
\par

\Equation{720\div20-8\textcolor{answerred}{<}720\div(20-8)}
\Equation{(75+25)\times4\textcolor{answerred}{>}75+25\times4}
\Equation{300\div15\times4\textcolor{answerred}{>}300\div(15\times4)}
\par\noindent

\Answer{
\QHeading{计算过程：第16题：比较大小}
\begin{enumerate}
\item 左边\(\begin{aligned}[t]&\phantom{={}}36\times9+36\\&=324+36\\&=360\end{aligned}\)，右边\(\begin{aligned}[t]&\phantom{={}}36\times(9+1)\\&=36\times10\\&=360\end{aligned}\)，填\(=\)。
\item 左边\(\begin{aligned}[t]&\phantom{={}}720\div20-8\\&=36-8\\&=28\end{aligned}\)，右边\(\begin{aligned}[t]&\phantom{={}}720\div(20-8)\\&=720\div12\\&=60\end{aligned}\)，填\(<\)。
\item 左边\(\begin{aligned}[t]&\phantom{={}}(75+25)\times4\\&=100\times4\\&=400\end{aligned}\)，右边\(\begin{aligned}[t]&\phantom{={}}75+25\times4\\&=75+100\\&=175\end{aligned}\)，填\(>\)。
\item 左边\(\begin{aligned}[t]&\phantom{={}}300\div15\times4\\&=20\times4\\&=80\end{aligned}\)，右边\(\begin{aligned}[t]&\phantom{={}}300\div(15\times4)\\&=300\div60\\&=5\end{aligned}\)，填\(>\)。
\end{enumerate}

}

\QHeading{第17题　解决问题。}

\Problem{①8名工人4小时做了480个零件。照这样计算，增加2名工人后，6小时可以做多少个零件？}
\par


\MethodTitle{方法一}
\Answer{\Formula{\begin{aligned}[t]&\phantom{={}}480\div4\div8\times(8+2)\times6\\&=480\div4\div8\times10\times6\\&=120\div8\times10\times6\\&=15\times10\times6\\&=150\times6\\&=900\text{（个）}\end{aligned}}\par 答：6小时可以做900个零件。}

\MethodTitle{方法二：先求原来8名工人6小时的产量}
\Answer{先把4小时的产量换算成6小时，再求新增2名工人的产量。}
\Answer{\textbf{原来8人6小时可做}\par \Formula{\begin{aligned}[t]&\phantom{={}}480\div 4\times 6\\
&=120\times 6\\
&=720\text{（个）}\end{aligned}}}
\Answer{\textbf{新增2人6小时可做}\par \Formula{\begin{aligned}[t]&\phantom{={}}720\div 8\times 2\\
&=90\times 2\\
&=180\text{（个）}\end{aligned}}}
\Answer{\textbf{合计}\par \Formula{\begin{aligned}[t]&\phantom{={}}720+180\\
&=900\text{（个）}\end{aligned}}}
\Answer{答：增加2名工人后，6小时可以做900个零件。}

\MethodTitle{方法三：先求一份工作量，再数总份数}
\Answer{1个人工作1小时算一份；先求原来有多少份，再求后来有多少份。}
\Answer{\textbf{原来的工作份数}\par \Formula{\begin{aligned}[t]&\phantom{={}}4\times 8\\
&=32\text{（份）}\end{aligned}}}
\Answer{\textbf{每份加工的零件数}\par \Formula{\begin{aligned}[t]&\phantom{={}}480\div 32\\
&=15\text{（个）}\end{aligned}}}
\Answer{\textbf{后来共有的份数}\par \Formula{\begin{aligned}[t]&\phantom{={}}6\times (8+2)\\
&=6\times 10\\
&=60\text{（份）}\end{aligned}}}
\Answer{\textbf{后来加工的零件数}\par \Formula{\begin{aligned}[t]&\phantom{={}}15\times 60\\
&=900\text{（个）}\end{aligned}}}
\Answer{答：6小时可以做900个零件。}
\Problem{②哥哥的身高是105厘米，比弟弟高20厘米。爸爸的身高正好是弟弟身高的2倍。爸爸比哥哥高多少厘米？}

\MethodTitle{方法一}
\Answer{\Formula{\begin{aligned}[t]&\phantom{={}}(105-20)\times2-105\\&=85\times2-105\\&=170-105\\&=65\text{（厘米）}\end{aligned}}\par 答：爸爸比哥哥高65厘米。}

\MethodTitle{方法二：直接比较爸爸与哥哥的差}
\Answer{爸爸是弟弟身高的2倍，也就是哥哥身高的2倍再减40厘米；再减一个哥哥的身高。}
\Answer{\Formula{\begin{aligned}[t]&\phantom{={}}105-20\times 2\\
&=105-40\\
&=65\text{（厘米）}\end{aligned}}}
\Answer{答：爸爸比哥哥高65厘米。}

